\subsection{Force--phase--delay audit loop: numerical differentiation, stability, and error budgets (audit)}
\label{app:force_phase_delay_audit}

This appendix records an audit-facing measurement/estimation loop that closes the interface chain
\[
\text{(force/response)}\ \Rightarrow\ \text{(phase)}\ \Rightarrow\ \text{(Wigner--Smith delay)}
\]
in a way that is reproducible and robust to the dominant numerical pitfalls (phase unwrapping and derivative noise amplification).
It is not a theorem-level statement about interacting QFT; it is a protocol for auditable inference once an $S(\omega)$-based platform is declared (Appendix~\ref{app:time_mass_delay} and Appendix~\ref{app:scattering_haag_ruelle_lsz_interface}).

\subsubsection{Objects and the numerical bottleneck}
\paragraph{Phase and delay.}
\InterfaceTag In one channel, $S(\omega)=\e^{\iu\delta(\omega)}$ and $\tau(\omega)=\dd\delta/\dd\omega$.
In multiple channels, the WS matrix and trace are
$Q(\omega)=-\iu S(\omega)^\dagger \dd S/\dd\omega$ and $\tau_{\mathrm{WS}}(\omega)=\Tr Q(\omega)$ (Appendix~\ref{app:time_mass_delay}, \eqref{eq:wigner_smith_omega}--\eqref{eq:tau_ws_trace_omega}).

\paragraph{Where instability enters.}
\AuditTag The map ``measured phase $\to$ delay'' is a \emph{frequency derivative}, and therefore amplifies noise.
Any claim about a delay perturbation $\Delta\tau(\omega)$ must be accompanied by a stated differentiation scheme and a stability sweep over the step size and smoothing.

\subsubsection{A minimal auditable pipeline (checklist)}
\begin{enumerate}
\item \textbf{Band selection.} Fix a frequency band $[\omega_{\min},\omega_{\max}]$ where scattering is approximately unitary to a stated tolerance (Appendix~\ref{app:time_mass_delay}, Remark~\ref{rem:wigner_smith_calibration_losses}).
\item \textbf{Phase extraction.} Extract $\delta(\omega)=\arg S(\omega)$ (one channel) or $\Theta(\omega)=\arg\det S(\omega)$ (multi-channel trace/logdet route), together with an explicit phase-unwrapping rule.
\item \textbf{Differentiation scheme.} Choose a finite-difference stencil for $\dd\delta/\dd\omega$ (or for $\dd\Theta/\dd\omega$) and declare the step size $\Delta\omega$.
\item \textbf{Smoothing (optional).} If smoothing is used, declare the kernel family and width; rerun all outputs across a bounded sweep of widths (no single ad-hoc choice).
\item \textbf{Stability sweep.} Recompute $\tau$ over a bounded family of $\Delta\omega$ and (if used) smoothing widths; report a stability envelope for the derived scalars (median, quantiles, max deviation).
\item \textbf{Cross-check.} Compare against an analytic benchmark when available (e.g.\ the one-channel Breit--Wigner profile; Appendix~\ref{app:time_mass_delay_reference}, \eqref{eq:breit_wigner_tau_omega}--\eqref{eq:tau_ws_peak}).
\end{enumerate}

\paragraph{Toy reproducibility check (optional).}
\AuditTag \paragraph{Toy audit summary (phase$\to$delay numerical stability).}
\AuditTag We generate a one-channel unitary Breit--Wigner model and estimate $\tau(\omega)=\mathrm{d}\delta/\mathrm{d}\omega$ from a noisy unwrapped phase under a bounded family of discretization choices (stride and smoothing window). The best median abs-log mismatch in the declared family is 0.022 (max 0.924) attained at stride=2, smoothing window w=21.
%
\begin{table}[t]
\centering
\small
\setlength{\tabcolsep}{10pt}
\renewcommand{\arraystretch}{1.15}
\begin{tabular}{lrrrr}
\toprule
variant & stride & smoothing window $w$ & median $|\log(\widehat\tau/\tau)|$ & max $|\log(\widehat\tau/\tau)|$ \\
\midrule
stride=1,w=1 & 1 & 1 & 0.456 & 6.259 \\
stride=1,w=7 & 1 & 7 & 0.125 & 4.838 \\
stride=1,w=21 & 1 & 21 & 0.041 & 2.162 \\
stride=2,w=1 & 2 & 1 & 0.290 & 5.005 \\
stride=2,w=7 & 2 & 7 & 0.074 & 3.344 \\
stride=2,w=21 & 2 & 21 & 0.022 & 0.924 \\
stride=4,w=1 & 4 & 1 & 0.195 & 5.852 \\
stride=4,w=7 & 4 & 7 & 0.038 & 0.902 \\
stride=4,w=21 & 4 & 21 & 0.024 & 0.588 \\
%
\end{tabular}
\caption{Toy audit for numerical stability of phase$\to$delay differentiation under a bounded family of discretization choices. Rows are generated by \texttt{scripts/exp\_force\_phase\_delay\_audit\_toy.py}.}
\label{tab:force_phase_delay_audit_toy}
\end{table}

\subsubsection{Error-budget template (finite differences)}
\paragraph{Derivative error decomposition.}
\AuditTag For a smooth phase function $\delta(\omega)$ sampled with additive phase noise $\eta(\omega)$, a central-difference estimator
\[
\widehat\tau(\omega):=\frac{\delta(\omega+\Delta\omega)-\delta(\omega-\Delta\omega)}{2\Delta\omega}
\]
admits the schematic decomposition
\[
\widehat\tau(\omega)-\tau(\omega)
=
O(\Delta\omega^2)\ +\ O\!\left(\frac{\sigma_\eta}{\Delta\omega}\right),
\]
where $\sigma_\eta$ is the typical phase-noise scale after unwrapping and any declared smoothing.
The first term is truncation error; the second is noise amplification.

\begin{table}[t]
\centering
\small
\setlength{\tabcolsep}{10pt}
\renewcommand{\arraystretch}{1.15}
\begin{tabular}{p{0.22\textwidth} p{0.33\textwidth} p{0.37\textwidth}}
\toprule
audit knob & bounded sweep (example) & what to report \\
\midrule
step size $\Delta\omega$ & $\Delta\omega\in\{\Delta\omega_0,2\Delta\omega_0,4\Delta\omega_0\}$ & envelope of $\tau(\omega)$ (max deviation; quantiles) \\
smoothing width (optional) & $h\in\{0,h_0,2h_0\}$ & bias vs variance tradeoff; note any drift of peaks \\
unitarity tolerance & $\|S^\dagger S-I\|$ threshold grid & whether $\tau_{\mathrm{WS}}$ is stable under the tolerance \\
phase unwrapping rule & bounded rule family & discontinuity count; phase-jump consistency \\
\bottomrule
\end{tabular}
\caption{Minimal stability-audit knobs for the phase$\to$delay derivative step. The purpose is to make the inference robust and auditable, not to optimize a single best-looking curve.}
\label{tab:force_phase_delay_audit_knobs}
\end{table}

\paragraph{Connection to protocol$\to$continuum error control.}
\AuditTag The error budgets in Appendix~\ref{app:protocol_to_continuum_error_control} address discretization and noise amplification in \emph{spatial} reconstruction pipelines (e.g.\ $\widehat\chi\mapsto \Delta_h\widehat\chi$).
The present appendix addresses the analogous amplification in \emph{frequency derivatives}.
Both are instances of the same audit discipline: whenever a derivative-like operator is applied, one must provide a bounded stability sweep and an explicit truncation/noise decomposition.
